\chapter{Data structures}
\begin{itemize}[nosep,leftmargin=*]
\item Public indices are 0-based; every range is $[l,r)$.
 Prefix arrays have n+1 entries. Internal Fenwick/DSU bookkeeping may differ.
\item T defaults to ll where applicable; arithmetic must fit the chosen type.
 Use sparse tables for static minima, Fenwick for additive queries,
 segment trees for general aggregates, lazy trees for range updates.
\end{itemize}
\section{Arrays and disjoint sets}
\kactlimport{ArrayBasics.h}
\kactlimport{UnionFind.h}
\kactlimport{UnionFindRollback.h}
\kactlimport{DynamicConnectivity.h}
\section{Range queries}
\kactlimport{Fenwick.h}
\kactlimport{Fenwick2D.h}
\kactlimport{SegmentTree.h}
\kactlimport{LazySegTree.h}
\kactlimport{SparseLazySegTree.h}
\kactlimport{RMQ.h}
\kactlimport{SqrtDecomposition.h}
\kactlimport{PrefixSum2D.h}
\kactlimport{Difference2D.h}
\section{Ordered and editable containers}
\kactlimport{OrderedSet.h}
\kactlimport{HashMap.h}
\kactlimport{Treap.h}
\par
